% The documents of the showcase's library, typeset with XeLaTeX by generate.ts. Written for the
% showcase: the people in them are fictional, and so are their words.
\documentclass[11pt]{article}
\usepackage[a4paper,margin=2.4cm]{geometry}
\usepackage{fontspec}
\usepackage{unicode-math}
\setmainfont{TeX Gyre Pagella}
\setmathfont{TeX Gyre Pagella Math}
\setmonofont{DejaVu Sans Mono}[Scale=0.84]
\usepackage{newunicodechar}
\newunicodechar{𝓞}{\ensuremath{\mathcal{O}}}
\usepackage{fancyvrb}
\usepackage{enumitem}
\usepackage{booktabs}
\usepackage{xcolor}
\usepackage[hidelinks]{hyperref}
\usepackage{titlesec}
\titleformat{\section}{\large\bfseries}{\thesection}{0.8em}{}
\titlespacing*{\section}{0pt}{1.2em}{0.4em}
\setlist{itemsep=0.15em,topsep=0.3em}
\setlength{\parindent}{0pt}
\setlength{\parskip}{0.45em}
\newcommand{\lean}[1]{\texttt{#1}}
\newcommand{\fictional}{\vfill{\footnotesize\color{gray}Written for the VisualProof showcase of
  flt-regular. Alice, Bob and Charlie are fictional, and so are their words; the mathematics and
  the Lean names are flt-regular's and Mathlib's.}}
\DefineVerbatimEnvironment{code}{Verbatim}{fontsize=\small,xleftmargin=1.2em,frame=leftline,framerule=0.6pt,rulecolor=\color{gray}}

\title{\bfseries Kummer's lemma in flt-regular:\\ reading notes}
\author{Charlie}
\date{8 September 2026}
\begin{document}
\maketitle
\thispagestyle{empty}

\section{The statement}

Let $p$ be an odd prime, $K$ a $p$-th cyclotomic field and $\mathcal{O}_K$ its ring of integers.
\textbf{Kummer's lemma.} If $p$ is regular and $u \in \mathcal{O}_K^\times$ is congruent to a
rational integer modulo $p$, then $u = v^p$ for a unit $v$. In the project:
\begin{code}
variable [IsCyclotomicExtension {p} ℚ K] (hp : p ≠ 2)
  (hreg : p.Coprime <| Fintype.card <| ClassGroup (𝓞 K))

theorem eq_pow_prime_of_unit_of_congruent (u : (𝓞 K)ˣ)
    (hcong : ∃ n : ℤ, (p : 𝓞 K) ∣ (u - n : 𝓞 K)) :
    ∃ v, u = v ^ p
\end{code}
Two things to check when approving it:
\begin{enumerate}
\item Regularity is stated for \lean{K}, as \lean{p.Coprime (Fintype.card (ClassGroup (𝓞 K)))},
  not as \lean{IsRegularPrime p}, which is about \lean{CyclotomicField p ℚ}. The two agree, since
  any two $p$-th cyclotomic fields over $\mathbb{Q}$ are isomorphic and the class number is
  invariant; the main theorem only ever applies the lemma to \lean{CyclotomicField p ℚ}.
\item The congruence is divisibility by $p$ in $\mathcal{O}_K$: $p \mid u - n$.
\end{enumerate}

\section{The second form}

\begin{code}
theorem exists_pow_eq_of_zeta_sub_one_pow_dvd_sub_one {u : (𝓞 K)ˣ}
    (hcong : (hζ.toInteger - 1 : 𝓞 K) ^ p ∣ (u : 𝓞 K) - 1) :
    ∃ v : K, v ^ p = u
\end{code}
Here the hypothesis is $(\zeta - 1)^p \mid u - 1$, and the conclusion is a $p$-th root in $K$, not
yet a unit. The first form is reduced to this one in the same file, through $u^{p-1}$: it is
congruent to $1$ modulo $p$, and from there to $1$ modulo $(\zeta - 1)^p$.

\section{How the proof goes}

Suppose $u$ is not a $p$-th power in $K$. Then $L = K(\sqrt[p]{u})$ is a cyclic extension of
degree $p$. The congruence $(\zeta - 1)^p \mid u - 1$ makes $L/K$ unramified: with
$\alpha^p = u$, the element $(1 - \alpha)/(\zeta - 1)$ is a root of
$\bigl(((\zeta - 1)X - 1)^p + u\bigr)/(\zeta - 1)^p$, a polynomial with integral coefficients,
separable modulo every prime (\lean{KummersLemma.isUnramified}, in \lean{KummersLemma.Field}). Hilbert's
Theorem~94 (\lean{dvd\_card\_classGroup\_of\_unramified\_isCyclic}) then says that $p$ divides the
class number of $K$, against regularity. Theorem~94 rests on Theorem~92 (\lean{Hilbert92}), on the
units of norm one in a cyclic extension of prime degree.

\section{For the review}

Approve the statements, not the proofs: Lean checked those. The class-group side
(\lean{Hilbert92}, \lean{dvd\_card\_classGroup\_of\_unramified\_isCyclic}) is worth reading for
what it states, since a weaker statement there would still compile.

\fictional
\end{document}
